% Chapter 4 worked example. Included by ch04_resources.tex.
\section{Worked Example: Ambient Documentation and Workforce Redesign}
\label{sec:comp-ch04-worked-example}

\subsection{Purpose and Status}
This worked example demonstrates how the methods in Chapter~4 can be
translated into a reviewable institutional decision. All numerical inputs are
synthetic unless a source is identified explicitly. They are intended to show the
logic of the analysis, not to report empirical findings or establish a universal
threshold. Local use requires replacement of the assumptions with current,
versioned evidence and approval through the relevant clinical, managerial,
economic, legal, technical, and governance processes.

An ambulatory network is considering an ambient documentation system to reduce after-hours work. The vendor reports substantial drafting-time savings, while clinicians are concerned about correction burden, privacy, coding changes, and pressure to convert every saved minute into additional appointments. 

\subsection{Decision Problem and Comparators}
\textbf{Decision problem.} Does the intervention create net workforce and organizational value, and what workflow, governance, and implementation conditions are required before scale?

The comparator set is deliberately broader than the existing pathway. This prevents
a new technology from receiving a lower evidentiary threshold than staffing,
workflow, shared-service, conventional digital, or no-change alternatives.

\begin{longtable}{@{}p{0.08\textwidth}p{0.84\textwidth}@{}}
\caption{Comparators for the Chapter 4 worked example}
\label{tab:comp-ch04-comparators}\\
\toprule
\textbf{Option} & \textbf{Comparator or strategy} \\
\midrule
\endfirsthead
\toprule
\textbf{Option} & \textbf{Comparator or strategy} \\
\midrule
\endhead
\bottomrule
\endlastfoot
1 & Current documentation workflow \\
2 & Template and workflow redesign without generative technology \\
3 & Additional transcription or administrative support \\
4 & Ambient documentation with opt-in human review \\
\end{longtable}

\subsection{Model and Decision Structure}
The analytical structure should follow the decision rather than the available
software. The team should map the causal pathway from intervention input to
information, human decision, action, outcome, resource consequence, and review.
The structure should also identify capacity constraints, uncertainty, subgroup
variation, implementation requirements, and events that would change the
recommendation.

A compact quantitative relationship used in this example is:
\begin{equation}
T_{\mathrm{net}}=T_{\mathrm{gross}}-T_{\mathrm{review}}-T_{\mathrm{correction}}-T_{\mathrm{coordination}}
\label{eq:comp-ch04-key-relationship}
\end{equation}
The equation is an organizing aid rather than a substitute for disaggregated evidence,
qualitative findings, or mandatory safeguards. Its variables and interpretation must
be defined in the local protocol.

\subsection{Synthetic Parameters and Evidence Sources}
Table~\ref{tab:comp-ch04-parameters} provides the base-case inputs. A
production analysis should add parameter distributions, correlations, structural
alternatives, data dates, system versions, and evidence-confidence ratings.

\begin{longtable}{@{}p{0.32\textwidth}p{0.22\textwidth}p{0.38\textwidth}@{}}
\caption{Synthetic parameters for the Chapter 4 worked example}
\label{tab:comp-ch04-parameters}\\
\toprule
\textbf{Parameter} & \textbf{Base value} & \textbf{Source or interpretation} \\
\midrule
\endfirsthead
\toprule
\textbf{Parameter} & \textbf{Base value} & \textbf{Source or interpretation} \\
\midrule
\endhead
\bottomrule
\endlastfoot
Clinicians in pilot & 60 & Pilot plan \\
Baseline documentation time per visit & 11.5 minutes & Time study \\
Vendor-reported gross saving & 5.0 minutes/visit & Vendor evidence \\
Observed review and correction time & 2.4 minutes/visit & Local usability test \\
Monthly subscription per clinician & EUR 420 & Proposed contract \\
Average visits per clinician per month & 310 & Scheduling system \\
Clinically consequential correction rate & 1.8\% & Pilot quality audit \\
\end{longtable}

\subsection{Step-by-Step Analysis}
The sequence below is intended to make the analysis reproducible and to ensure
that evidence, economics, governance, and implementation develop together.

\begin{longtable}{@{}p{0.07\textwidth}p{0.55\textwidth}p{0.30\textwidth}@{}}
\caption{Analytical sequence}
\label{tab:comp-ch04-analysis-steps}\\
\toprule
\textbf{Step} & \textbf{Required work} & \textbf{Decision record} \\
\midrule
\endfirsthead
\toprule
\textbf{Step} & \textbf{Required work} & \textbf{Decision record} \\
\midrule
\endhead
\bottomrule
\endlastfoot
1 & Map the full documentation workflow and distinguish drafting assistance from coding or downstream recommendations. & Evidence and responsible owner to be recorded \\
2 & Measure gross time, review, correction, escalation, and coordination by professional role. & Evidence and responsible owner to be recorded \\
3 & Assess whether net time can be pooled or scheduled as usable capacity. & Evidence and responsible owner to be recorded \\
4 & Evaluate quality, privacy, consent, patient interaction, and after-hours work alongside throughput. & Evidence and responsible owner to be recorded \\
5 & Use CFIR, NASSS-CAT, or another selected implementation resource for a defined purpose rather than as a generic checklist. & Evidence and responsible owner to be recorded \\
6 & Define benefit owners, renewal thresholds, and controls for model, prompt, and template updates. & Evidence and responsible owner to be recorded \\
\end{longtable}

\subsection{Illustrative Outputs}
The outputs in Table~\ref{tab:comp-ch04-outputs} show how technical,
clinical, operational, economic, equity, and governance findings should be kept
visible rather than compressed into a single favourable or unfavourable score.

\begin{longtable}{@{}p{0.30\textwidth}p{0.24\textwidth}p{0.38\textwidth}@{}}
\caption{Illustrative model and decision outputs}
\label{tab:comp-ch04-outputs}\\
\toprule
\textbf{Output} & \textbf{Result} & \textbf{Interpretation} \\
\midrule
\endfirsthead
\toprule
\textbf{Output} & \textbf{Result} & \textbf{Interpretation} \\
\midrule
\endhead
\bottomrule
\endlastfoot
Gross time saved & 5.0 minutes/visit & Vendor claim \\
Net time saved after review & 2.6 minutes/visit & Local pilot estimate \\
Potential monthly net time per clinician & 13.4 hours & Before accounting for downtime and meetings \\
Cash-releasing saving & Not established & No payroll reduction demonstrated \\
Main realized value candidate & Reduced after-hours work and improved patient interaction & Requires longitudinal evidence \\
\end{longtable}

\subsection{Sensitivity, Uncertainty, and Subgroups}
At minimum, the completed analysis should vary the parameters most likely to
change the decision, test at least one credible alternative model structure, and
report subgroup results separately where performance, access, response, burden,
or opportunity cost may differ. Deterministic analysis should identify thresholds;
probabilistic analysis should be used where credible parameter distributions and
correlations are available. Scenario analysis is preferable when structural or
future uncertainty cannot be represented defensibly by one probability model.

The record should distinguish uncertainty that can be reduced through evidence,
uncertainty that can be managed through safeguards, and uncertainty that remains
too consequential to accept. Missing evidence should not be represented as an
average or neutral value without justification.

\subsection{Interpretation and Recommendation}
Continue a limited opt-in pilot and do not treat gross time saving as a cash benefit. Scale only if net time benefit, documentation quality, patient acceptability, privacy, and workforce experience remain favourable and no adverse coding or downstream-utilization signal emerges.

The recommendation is conditional rather than permanent. It applies only to the
specified population, setting, intervention configuration, evidence cut-off, and
version. The following conditions should be incorporated into the approval,
contract, implementation, monitoring, or renewal record:

\begin{itemize}
 \item documentation assistance remains separate from coding and order automation.
 \item clinically consequential content requires meaningful professional confirmation.
 \item specialty-specific correction patterns and subgroup effects are monitored.
 \item released time has a documented and acceptable use.
 \item renewal is linked to sustained net benefit after the novelty period.
\end{itemize}

\subsection{Decision Statement}
\begin{quote}
For the specified decision problem and comparator set, the institution should
\emph{[proceed / proceed with conditions / redesign / pilot / restrict / defer /
replace / retire]}. The decision applies to \emph{[population, setting, version,
and evidence period]}, is owned by \emph{[accountable authority]}, and will be
reviewed on \emph{[date]} or earlier if \emph{[trigger events]} occur. The
evidence, assumptions, unresolved uncertainty, mandatory safeguards, monitoring
thresholds, and exit arrangements are recorded in the accompanying dossier.
\end{quote}
